\section{Kiểm thử hệ thống}
\subsection{Mục tiêu kiểm thử}
Sau khi hoàn thiện quá trình phân tích, thiết kế và hiện thực hệ thống,
các chức năng chính của CityFoodTour được tiến hành kiểm thử nhằm đảm bảo
hệ thống hoạt động đúng theo yêu cầu đã đề ra.

Quá trình kiểm thử tập trung vào việc xác minh tính đúng đắn của các chức năng,
khả năng xử lý các luồng nghiệp vụ chính và sự tương tác giữa các thành phần
trong hệ thống. Các chức năng được kiểm thử bao gồm quản lý người dùng,
tìm kiếm quán ăn, quản lý food tour, gợi ý quán ăn, đánh giá và các chức năng
quản trị hệ thống.

\subsection{Môi trường kiểm thử}
Hệ thống được kiểm thử trong môi trường phát triển với các thành phần:

\begin{itemize}

    \item \textbf{Frontend:}
    React kết hợp TypeScript, được chạy trên môi trường trình duyệt web.

    \item \textbf{Backend:}
    Node.js và Express.js chịu trách nhiệm xử lý các API và nghiệp vụ hệ thống.

    \item \textbf{Database:}
    MongoDB được sử dụng để lưu trữ thông tin người dùng, quán ăn,
    food tour và các dữ liệu liên quan.

    \item \textbf{Công cụ kiểm thử:}
    Postman được sử dụng để kiểm tra các API,
    kết hợp với việc kiểm thử trực tiếp trên giao diện người dùng.

\end{itemize}

\subsection{Kiểm thử chức năng}

Các chức năng chính của hệ thống được kiểm thử dựa trên các kịch bản sử dụng
thực tế nhằm đảm bảo hệ thống đáp ứng đúng các yêu cầu đã đề ra.

\begin{longtable}{|c|p{3cm}|p{7cm}|c|}
\caption{Danh sách kiểm thử chức năng}
\label{tab:testcase}\\

\hline
STT & Chức năng & Kết quả mong đợi & Kết quả \\
\hline

1 &
Đăng ký tài khoản &
Người dùng có thể tạo tài khoản mới với thông tin hợp lệ.
Hệ thống thông báo lỗi khi thông tin đăng ký không hợp lệ.
&
Đạt
\\
\hline

2 &
Đăng nhập &
Người dùng đăng nhập thành công với thông tin chính xác.
Hệ thống từ chối khi sai thông tin đăng nhập.
&
Đạt
\\
\hline

3 &
Cập nhật hồ sơ người dùng &
Người dùng có thể cập nhật thông tin cá nhân và sở thích ẩm thực.
&
Đạt
\\
\hline

4 &
Tìm kiếm quán ăn &
Hệ thống trả về danh sách quán ăn dựa trên từ khóa hoặc bộ lọc.
&
Đạt
\\
\hline

5 &
Xem chi tiết quán ăn &
Người dùng có thể xem thông tin chi tiết của một quán ăn.
&
Đạt
\\
\hline

6 &
Thêm quán ăn yêu thích &
Người dùng có thể lưu và quản lý danh sách quán ăn yêu thích.
&
Đạt
\\
\hline

7 &
Tạo food tour &
Người dùng có thể tạo mới một hành trình ăn uống.
&
Đạt
\\
\hline

8 &
Thêm địa điểm vào food tour &
Người dùng có thể thêm các quán ăn vào hành trình đã tạo.
&
Đạt
\\
\hline

9 &
Chỉnh sửa food tour &
Người dùng có thể cập nhật thông tin hoặc danh sách địa điểm trong food tour.
&
Đạt
\\
\hline

10 &
Gợi ý quán ăn &
Hệ thống hiển thị danh sách quán ăn phù hợp dựa trên thông tin người dùng.
&
Đạt
\\
\hline

11 &
Đánh giá quán ăn &
Người dùng có thể gửi đánh giá cho quán ăn đã sử dụng.
&
Đạt
\\
\hline

12 &
Chia sẻ food tour &
Người dùng có thể chuyển food tour sang trạng thái công khai để chia sẻ.
&
Đạt
\\
\hline

13 &
Đăng tải nhà hàng &
Người dùng có quyền quản lý nhà hàng có thể gửi yêu cầu đăng tải nhà hàng.
&
Chưa kiểm thử
\\
\hline

14 &
Duyệt nhà hàng &
Quản trị viên có thể xem, phê duyệt hoặc từ chối yêu cầu đăng tải nhà hàng.
&
Chưa kiểm thử
\\
\hline

15 &
Quản lý tài khoản quản trị &
Quản trị viên có thể quản lý thông tin tài khoản người dùng.
&
Đạt
\\
\hline

\end{longtable}

\subsection{Kiểm thử các module chính}
\subsubsection{Kiểm thử module quản lý người dùng}
\subsubsection{Kiểm thử module tìm kiếm quán ăn}
\subsubsection{Kiểm thử module Food Tour}
\subsubsection{Kiểm thử module gợi ý quán ăn}
\section{Đánh giá kết quả kiểm thử}

Thông qua quá trình kiểm thử, các chức năng chính của hệ thống đều hoạt động
theo đúng yêu cầu thiết kế. Các luồng xử lý quan trọng như đăng nhập,
tìm kiếm quán ăn, xây dựng food tour, quản lý thông tin và gợi ý quán ăn
đều cho kết quả phù hợp.

Kết quả kiểm thử cho thấy kiến trúc hệ thống và sự phân tách giữa các tầng
frontend, backend và database hoạt động ổn định. Tuy nhiên, hệ thống vẫn còn
một số giới hạn như dữ liệu thử nghiệm chưa bao phủ toàn bộ các trường hợp
thực tế và cần tiếp tục được đánh giá với tập dữ liệu lớn hơn trong tương lai.
